\documentclass{article}
\usepackage[UTF8]{ctex}
\usepackage{amsmath,amssymb,amsfonts}
\usepackage{enumitem}
\usepackage{geometry}
\geometry{a4paper, margin=1in}
% === 错题本排版优化（不参与题目标识） ===
\setlength{\parindent}{0pt}
\setlength{\parskip}{0.35em plus 0.10em minus 0.08em}
\setlength{\emergencystretch}{3em}
\setlist{itemsep=0.25em, topsep=0.40em, parsep=0pt}
\allowdisplaybreaks[2]
% === 错题本排版优化结束 ===
\begin{document}
\section*{哈尔滨工业大学2024年数学分析真题}
\section*{81．哈尔滨工业大学2024年数学分析试卷}

\subsection{题目1}
1．正数列 \(\left\{a_{n}\right\},\left\{b_{n}\right\}\) 发散，则 \(\left\{a_{n}+b_{n}\right\}\) 发散．

\subsection{题目2}
2．\(f(x)\) 在 \((-1,1)\) 上可微，\(f^{\prime}(0)>0\) ，则存在 \(\varepsilon>0\) 满足 \(f(x)\) 在邻域 \(O(0, \varepsilon) \subset(-1,1)\) 上单调。

\subsection{题目3}
3．若 \(\int_{a}^{b}(f(x))^{2024} \mathrm{~d} x=0\) ，则 \(f(x) \equiv 0, x \in[a, b]\) ．

\subsection{题目4}
4．设非负项级数 \(\sum_{n=1}^{\infty} a_{n}, \sum_{n=1}^{\infty} b_{n}\) 收敛，则 \(\sum_{n=1}^{\infty} \min \left\{a_{n}, b_{n}\right\}\) 收敛。

\subsection{题目5}
5．设 \(f(x, y)\) 沿任意方向趋于 \((0,0)\) 的方向极限都存在且相等，则重极限 \(\lim _{(x, y) \rightarrow(0,0)} f(x, y)\) 一定存在．

二．设 \(m, n \geq 1\) ，求 \(\lim _{x \rightarrow 1}\left(\frac{m}{1-x^{m}}-\frac{n}{1-x^{n}}\right)\) ．

三．\(f(x)\) 是定义在 \((a, b)\) 上的函数，对于任意收敛数列 \(\left\{x_{n}\right\} \subset(a, b),\left\{f\left(x_{n}\right)\right\}\) 收敛．证明：

\begin{itemize}
  \item[（1）] \(f(x) \in C(a, b)\) ．
  \item[（2）] \(\lim _{x \rightarrow a^{+}} f(x), \lim _{x \rightarrow b^{-}} f(x)\) 存在．
  \item[（3）] \(f(x)\) 在 \((a, b)\) 上一致连续．
\end{itemize}

四．设 \(f(x)\) 在 \(\mathbb{R}\) 上二阶可微，对任意的 \(x \in \mathbb{R},|f(x)| \leq M_{0},\left|f^{\prime \prime}(x)\right| \leq M_{2}\) ，

\begin{itemize}
  \item[（1）] 写出 \(\forall x \in \mathbb{R}, \forall h>0, f(x+h), f(x-h)\) 关于 \(h\) 带有 Lagrange 余项的 Taylor 公式。
  \item[（2）] 证明：\(\forall x \in \mathbb{R}, \forall h>0,|f^{\prime}(x)| \leq \frac{M_{0}}{h}+\frac{h}{2} M_{2}\) ．
  \item[（3）] 证明：\(\forall x \in \mathbb{R},|f^{\prime}(x)| \leq \sqrt{2 M_{0} M_{2}}\) ．
\end{itemize}

五、解答如下问题：

\begin{itemize}
  \item[（1）] 证明： \(\lim _{b \rightarrow 1^{-}} \int_{0}^{b} \frac{\cos x}{\sqrt{1-x^{2}}} \mathrm{~d} x\) 存在．
  \item[（2）] 证明： \(\lim _{b \rightarrow 1^{-}} \int_{0}^{b} \frac{\cos x}{\sqrt{1-x^{2}}} \mathrm{~d} x>1\) ．
\end{itemize}

六。设 \(\left\{a_{n}\right\}\) 是单调递增的正数列，证明：

\begin{itemize}
  \item[（1）] 若 \(\left\{a_{n}\right\}\) 有界，则 \(\sum_{n=1}^{\infty}\left(1-\frac{a_{n}}{a_{n+1}}\right)\) 收敛。
  \item[（2）] 若 \(\left\{a_{n}\right\}\) 无界，则 \(\sum_{n=1}^{\infty}\left(1-\frac{a_{n}}{a_{n+1}}\right)\) 发散。
\end{itemize}

七．设 \(\left\{f_{n}(x)\right\}=\left\{x n^{\alpha} e^{-n x}\right\}, n=1,2, \cdots\) ．

\begin{itemize}
  \item[（1）] 当 \(\alpha\) 为何值时，\(\left\{f_{n}(x)\right\}\) 在［0，1］上收敛．
  \item[（2）] 当 \(\alpha\) 为何值时，\(\left\{f_{n}(x)\right\}\) 在［0，1］上一致收敛。
  \item[（3）] 讨论 \(\lim _{n \rightarrow \infty} \int_{0}^{1} f_{n}(x) \mathrm{d} x\) 和 \(\int_{0}^{1} \lim _{n \rightarrow \infty} f_{n}(x) \mathrm{d} x\) 的关系．
\end{itemize}

八．设 \(\varphi(x, y)\) 在 \((0,0)\) 的某一邻域上连续，\(f(x, y)=|x-y| \varphi(x, y)\) ．

\begin{itemize}
  \item[（1）] 当 \(\varphi(x, y)\) 满足什么条件时，\(f_{x}^{\prime}(0,0), f_{y}^{\prime}(0,0)\) 存在？
  \item[（2）] 当 \(\varphi(x, y)\) 满足什么条件时，\(f(x, y)\) 在（0，0）处可微？
\end{itemize}

九．求 \(\lim _{n \rightarrow \infty} \frac{1}{n^{4}} \iiint_{x \leq n}[r] \mathrm{d} x \mathrm{~d} y \mathrm{~d} z\) ，其中 \(r=\sqrt{x^{2}+y^{2}+z^{2}}\) ．

\subsection{题目14}
原题附加题：求 \(\iint_{S} x^{2} \mathrm{~d} y \mathrm{~d} z+y^{2} \mathrm{~d} z \mathrm{~d} x+z^{2} \mathrm{~d} x \mathrm{~d} y\) ，其中 \(S\) 是 \(z=x^{2}+y^{2}\) 从 \(z=0\) 到 \(z=1\) 的部分，取下侧。
\end{document}
